\section{Introduction}
\label{sec:introduction}

\lettrine{A}{ccurate} characterization of wing aerodynamics requires not only integrated force and moment coefficients but also the surface pressure distribution, which provides spatially resolved information on aerodynamic loading, shock-wave behavior, and other local flow features. Reynolds-averaged Navier--Stokes (RANS) computational fluid dynamics (CFD) remains one of the primary tools for obtaining such high-fidelity aerodynamic information. However, aerodynamic design typically requires repeated evaluations over variations in geometry and operating conditions, causing the computational cost of CFD to increase rapidly in multipoint analysis, parametric studies, and aerodynamic optimization. This issue becomes particularly pronounced for transonic wings, where shock waves, shock--boundary-layer interactions, and possible flow separation introduce strong nonlinearities into the aerodynamic response. Consequently, developing rapid predictive models that retain the spatial information of high-fidelity CFD while substantially reducing the evaluation cost has become an important topic in data-driven aerodynamics \cite{Brunton2020MLFluid}.

Deep learning provides an effective means of constructing aerodynamic surrogate models by learning nonlinear mappings between geometry, flow conditions, and distributed aerodynamic quantities. For two-dimensional airfoils, convolutional and encoder--decoder neural networks have been successfully used to predict pressure and velocity fields over previously unseen geometries \cite{Bhatnagar2019CNNFlow,Thuerey2020DeepRANS}. Subsequent studies have demonstrated accurate prediction of airfoil surface pressure distributions and transonic flow fields involving shock waves and flow separation \cite{Hui2020FastPressure,Duru2022TransonicFlow}. Compared with surrogate models that predict only integrated aerodynamic coefficients, distributed-field prediction preserves local aerodynamic information and therefore provides a more informative basis for load evaluation, flow-feature analysis, and aerodynamic design. More recently, data-driven models have also been extended to three-dimensional configurations, including aerodynamic surface-pressure prediction on general three-dimensional geometries and reduced-order prediction of transonic wing flow fields \cite{Shen2023General3D,Immordino2024TransonicDL}. These developments demonstrate the potential of deep learning for rapid three-dimensional aerodynamic prediction, while also exposing a fundamental difficulty: generating sufficiently diverse high-fidelity three-dimensional training data is considerably more expensive than generating corresponding two-dimensional data.

One possible route to reducing this data requirement is to exploit the physical relationship between lower-dimensional sectional aerodynamics and finite-wing flows. Classical sweep corrections and quasi-three-dimensional approaches seek to approximate three-dimensional wing aerodynamics through appropriately transformed two-dimensional or infinite-swept-wing solutions. In particular, 2.75D approaches incorporate sweep- and taper-related effects into sectional analyses and have been successfully employed for transonic wing analysis and design at a computational cost substantially lower than that of full three-dimensional optimization \cite{Zhao2016SupercriticalWing,Xu2022Improved275D}. A similar principle has recently been introduced into data-driven modeling. Li et al. transferred aerodynamic knowledge learned from two-dimensional supercritical airfoils to three-dimensional transonic swept wings \cite{Li2023AirfoilToWing}, while Yang et al. further embedded sweep theory into an airfoil-to-wing transfer-learning framework \cite{Yang2025PhysicsEmbeddedTransfer}. These studies indicate that lower-dimensional aerodynamic information contains transferable pressure characteristics and flow features that can provide useful priors for three-dimensional prediction.

Nevertheless, the relationship between sectional aerodynamics and a finite wing is inherently nonuniform along the span. Section-based or conical-flow approximations are most applicable where the local flow remains approximately quasi-two-dimensional, whereas their assumptions progressively deteriorate in regions dominated by finite-wing effects. In fact, the applicability of the 2.75D transformation has been reported to be limited in regions with strong three-dimensional effects, particularly near the wing root and wing tip \cite{Xu2022Improved275D}. Experimental investigations of finite transonic wings have likewise shown that sweep, spanwise flow, and tip-vortex effects can substantially modify the local flow relative to the corresponding two-dimensional configuration \cite{DAguanno2022FiniteWing}. Under transonic conditions, spanwise variations in shock location and strength further enhance this spatial nonuniformity. Consequently, a sectional pressure model may reproduce the dominant chordwise pressure characteristics over much of the wing while still exhibiting systematic local discrepancies in the root and tip regions. Accurately representing these regional three-dimensional effects is therefore essential if lower-dimensional aerodynamic knowledge is to be used as a reliable prior for full finite-wing pressure prediction.

The limited availability of three-dimensional training data provides a second challenge. Transfer learning and multifidelity learning have therefore been increasingly introduced into aerodynamic modeling to reuse information obtained from related configurations or lower-cost data. Multifidelity neural-network models have demonstrated that inexpensive aerodynamic data can be combined with a smaller number of high-fidelity samples to improve aerodynamic surrogate modeling \cite{Zhang2021MultifidelityDNN,Li2022Multifidelity,Leng2024VariableFidelity}. For flow-field prediction, Wang et al. employed transfer learning to improve airfoil-flow prediction when only a small target-domain dataset was available \cite{Wang2023DeepTransfer}. At the wing level, spatial dependencies learned from available pressure information have been transferred between spanwise sections \cite{Qu2022SpatialTransfer}, while transferable neural-network architectures have been used to reduce the number of samples required when extending aerodynamic models across swept-wing configurations \cite{Yang2024TransferableWing}. Physics-embedded airfoil-to-wing transfer further demonstrates that explicitly incorporating sectional aerodynamic knowledge can substantially reduce the amount of three-dimensional wing data required for model construction \cite{Yang2025PhysicsEmbeddedTransfer}. However, transferred knowledge remains effective only insofar as the source-domain representation captures the relevant physics of the target flow; discrepancies associated with finite-span and strongly three-dimensional local effects cannot generally be eliminated by transferring sectional information alone.

The existing literature therefore reveals a central tradeoff in small-sample finite-wing prediction. Lower-dimensional aerodynamic data are relatively inexpensive to generate and contain reusable information about the fundamental chordwise pressure response, but they cannot fully represent spatially localized three-dimensional effects. High-fidelity finite-wing data retain these effects directly, but obtaining sufficient samples to cover the combined geometry--condition space is computationally expensive. Thus, the key issue is not simply how to transfer two-dimensional knowledge to a three-dimensional model, but how to distinguish aerodynamic information that is broadly reusable from that associated with specific finite-wing regions, and how to combine both forms of knowledge with a limited amount of target-domain data.

To address this issue, the present study develops a hierarchical aerodynamic-knowledge-driven framework for predicting the surface pressure distribution and aerodynamic force coefficients of finite symmetric swept wings under transonic conditions. The framework first employs lower-dimensional aerodynamic data to learn the fundamental sectional pressure characteristics. The pressure difference between finite and corresponding infinite swept wings is then used to identify localized three-dimensional corrections associated with the wing-root and wing-tip regions, thereby introducing regional aerodynamic knowledge beyond the sectional representation. These knowledge components are combined to construct a physics-informed pressure prior for the target finite wing. Finally, a limited number of high-fidelity target-wing samples are used to learn the residual discrepancy between this prior and the actual three-dimensional pressure field. In this manner, direct learning of the complete three-dimensional aerodynamic response is reformulated as a hierarchical correction process from reusable sectional knowledge to regional three-dimensional effects and, ultimately, to target-specific residuals. 
The predicted surface pressure distributions are further used, together with geometric and flow-condition information, to construct lightweight regression models for the integrated aerodynamic force and moment coefficients. This enables the assessment of whether the benefits of hierarchical aerodynamic knowledge transfer extend from distributed pressure prediction to integrated aerodynamic responses under limited target-domain training data.